\section{Introduction}
\label{sec:intro}

Large language models are trained to be both versatile next-token predictors and helpful conversational agents. Pre-training on diverse corpora creates pattern completers that can learn from context at inference time through in-context learning \citep{dongSurveyIncontextLearning2023, riechers_next-token_2025} and simulate a variety of personas \citep{nostalgebraist2025void}. Post-training, through supervised finetuning, RLHF \citep{baiTrainingHelpfulHarmless2022}, DPO \citep{rafailov2023direct}, and related techniques---then attempts to channel this flexibility into a stable, helpful assistant persona that is consistently presented to the user \citep{luAssistantAxisSituating2026, marks2026psm}, regardless of context. The same ``porosity'' to context that allows the model to adapt comes with its risks. It can cause a model to infer and adopt a broadly misaligned persona from a few examples \citep{afonin_emergent_2026}, and it can be exploited by jailbreaking attacks that construct conversational histories that portray the model as having already complied with harmful requests \citep{anil_many-shot_2024}, or that gradually escalate from benign exchanges to harmful ones \citep{russinovich2025great}. A key goal of post-training is therefore to inoculate models against such context-driven behavioral updating that could lead to misalignment or other undesirable behavior, while still allowing them to track factual and conversational content in service of being helpful assistants.

Instruction-following robustness has been evaluated along numerous dimensions \citep{pyatkin_generalizing_2025, qin_infobench_2024, zhou_instruction-following_2023, wen_benchmarking_2024}, but existing benchmarks do not evaluate this robustness in the presence of conflicting evidence about the model's behavior.
In this paper, we introduce a controlled paradigm for studying this \emph{induction-instruction conflict}.
The model receives an explicit instruction specifying target behavior~$T$, followed by $N$ hardcoded assistant turns producing a competing behavior~$P$.
As $N$ increases, inductive pressure grows, placing it in direct conflict with the instruction.
We evaluate 13 models across 16 instruction types and ask: at what~$N$ does the model abandon its instruction? How does this depend on the model, and the type and content of the instruction?
We make no normative claim about which behavior is correct; our aim is to characterize \emph{where} this transition falls and what governs it.

We find that the transition from instruction-following to pattern-following is universal but far from uniform.
Average instruction-following rates range from 1\% to 99\% across models, largely uncorrelated with standard capability benchmarks such as GPQA \citep{rein_gpqa_2023}.
Behavior is modulated by instruction content: models resist induction longer when instructions align with trained value priors. Enabling reasoning improves performance, but does not guarantee instruction adherence.

Our key contributions are: (1) a minimal, controlled paradigm for quantifying instruction-induction conflict across models and conditions; (2) a systematic evaluation of 13 models showing that the transition is universal, model-dependent, and modulated by the type and content of the instruction; (3) evidence that post-training method and reasoning capability both influence robustness, and that output diversity, not semantic engagement with the input, is the primary factor distinguishing robust from fragile conditions. The remainder of the paper is organized as follows. Section \ref{sec:methods} describes the experimental paradigm, conditions, models, and evaluation protocols. Section \ref{sec:results} presents results, Section \ref{sec:discussion} discusses implications and limitations, and Section \ref{sec:conclusion} concludes. Related work is discussed in Section~\ref{app:related}.
